% GateTruth v1.0 paper (renamed from SiliconBench per ADR-0008; pivot to an external
% mutation-testing audit as the headline contribution, per the same ADR).
% NOTE: This was assembled but NOT compiled locally (no TeX toolchain in the build
% environment). Compile with:  latexmk -pdf main.tex   (runs bibtex automatically),
% or on Overleaf. Requires the auto-generated table fragments under data/build/*.tex
% (regenerate with: python paper/data/generate_tables.py).
% STATUS: Full pivot (ADR-0008) complete. External Audit Findings (sec:audit-results) now
% reports real, independently-reproduced numbers from vault tickets SB-109..SB-112
% (Architect-verified from fresh, independent RTLLM clones -- see vault ticket evidence).
\documentclass[11pt]{article}

\usepackage[utf8]{inputenc}
\usepackage[T1]{fontenc}
\usepackage[margin=1in]{geometry}
\usepackage{amsmath}
\usepackage{booktabs}
\usepackage{graphicx}
\usepackage{microtype}
\usepackage{enumitem}
\usepackage{xcolor}
\usepackage[hidelinks]{hyperref}

\title{GateTruth: Auditing the Rigor of RTL Design Benchmarks via Mutation Testing}
\author{Meet Bhadra \\ Independent Researcher}
\date{\today}

\begin{document}
\maketitle

\begin{abstract}
Benchmarks for evaluating large language models on register-transfer-level (RTL) hardware design
have proliferated rapidly, yet none of them has applied mutation testing --- an established
hardware-verification technique for quantifying testbench quality --- to ask whether its own
testbenches are trustworthy. A testbench that never fails is not evidence of a correct design; it
may simply never stimulate the logic that is actually broken. We introduce GateTruth, a
mutation-testing engine and methodology for auditing the rigor of RTL benchmark testbenches: we
inject a deterministic, seeded set of semantic mutants into a reference design and measure what
fraction the testbench catches. A benchmark whose testbenches pass a deliberately broken design is
not measuring what it claims to measure. We validate the methodology first against our own
68-task, dual-track reference suite --- 60 specification-to-RTL generation tasks and 8
agentic-repair tasks, each scored through a pinned, deterministic synthesis-to-timing flow with
correctness enforced as a strict gate rather than a weighted term --- certifying that every one of
the 60 Track A testbenches kills at least 95\% of injected mutants under sequential, reproducible
execution. We then point the same engine, unmodified, at RTLLM v2.0, a widely adopted external
RTL-generation benchmark, and report the measured mutation-kill rates of its public testbenches
against its own reference implementations: of the 44 designs whose reference (as shipped, or after
a documented compile-time build-convention alias) passes baseline validation, 73\% fall below the
95\% floor our own suite is held to, and three score a 0\%
kill rate outright. One of the six designs excluded from this count is excluded for a different
reason entirely: its own shipped golden reference fails its own shipped testbench outright,
independent of mutation. Where a comparable audit of NVIDIA's CVDP benchmark is
not possible --- its public release deliberately withholds reference solutions to prevent
contamination, which also removes the golden RTL that mutation testing requires --- we report that
gap as a structural finding in its own right rather than silently omitting the benchmark. We detail
the audit methodology, the reference-suite design used to validate it, and the resulting kill-rate
measurements, and argue that mutation-kill certification should become a standard reporting
requirement for RTL-generation benchmarks generally.
\end{abstract}

\section{Introduction}
\label{sec:intro}

The evaluation of large language models on hardware description languages has converged on a
standard recipe: give the model a natural-language specification, ask it to generate Verilog or
SystemVerilog, and score pass@k against a testbench. This recipe answers one question --- did the
generated code pass the tests --- and treats the testbench itself as ground truth. That assumption
is rarely checked. A testbench with weak coverage will pass a broken design as readily as it passes
a correct one, and a benchmark built on such testbenches silently overstates every model's
performance, in a way that pass@k alone cannot reveal: the number looks identical whether the
testbench is rigorous or vacuous.

Mutation testing is the standard software-engineering answer to exactly this problem: inject small,
deterministic, semantically meaningful faults into a reference implementation and measure what
fraction of them a test suite detects (``kills''). A testbench that fails to catch an injected
sign flip, an inverted comparator, or a dropped reset condition is a testbench that would also fail
to catch the analogous mistake in a model-generated design. Despite its maturity elsewhere in
software engineering, mutation-kill certification is, to our knowledge, not reported by any
existing public RTL-generation benchmark for LLMs.

This gap is addressable directly: build a mutation engine, certify it against a reference suite
authored specifically to validate the methodology, and then point the same engine, unmodified, at
benchmarks the field already relies on. That is GateTruth's contribution. The reference suite ---
60 specification-to-RTL tasks across three complexity tiers, plus 8 agentic-repair tasks that
evaluate whether a tool-augmented agent can improve an existing design's physical characteristics or
extend it with a required new property (e.g., latch removal, CDC safety) under a strict correctness
gate --- exists to give the mutation engine, the correctness-as-gate
scoring discipline, and the contamination-defense apparatus a fully controlled environment in which
to be built and certified before being trusted against external artifacts. The physical
power/performance/area (PPA) scoring built into that reference suite is real and independently
useful (Section~\ref{sec:results} reports it), but it is not this paper's headline: the headline is
what the same rigor-checking apparatus finds when pointed outward.

This paper makes three contributions. First, a mutation-testing audit methodology for RTL
benchmarks --- a deterministic, seeded, sequential-execution protocol for measuring testbench
kill rates against external reference implementations, together with an explicit policy for
designs whose testbenches cannot be run at all (reported, never silently dropped or scored as a
misleading zero). Second, an openly published, dual-track reference implementation of that
methodology: 60 specification-to-RTL tasks and 8 agentic-repair tasks, each scored through a
pinned, deterministic synthesis-to-timing flow, with every one of the suite's 60 Track A
testbenches certified above a 95\% mutation-kill floor under the audit protocol before any external
use --- because a tool that audits other benchmarks' rigor must first demonstrate its own
(Track B's 8 testbenches are not yet certified under this protocol; see Section~\ref{sec:limitations}).
Third, the audit results
themselves (Section~\ref{sec:audit-results}): of RTLLM v2.0's 50 public designs, 44 pass baseline
validation and are audited, and 73\% of those fall below the 95\% mutation-kill floor our own
suite is certified against, with three of the audited designs at a 0\% kill rate. Among the six
designs that do not pass baseline validation and are therefore excluded from the 44, one is
excluded for a notable reason distinct from the rest: its own shipped golden reference fails its
own shipped testbench outright --- plus a documented structural finding,
rather than a silent gap, for why an equivalent measurement against NVIDIA's CVDP benchmark is
not currently possible from its public release.

\section{GateTruth Design}
\label{sec:design}

GateTruth evaluates large language model (LLM) and agent-generated register-transfer-level (RTL)
designs along two distinct axes frequently conflated or omitted by prior evaluation frameworks:
functional correctness and physical viability---specifically, optimization for area, timing, and
power relative to a competent human-authored baseline. This design serves two purposes in this
paper: it is a useful benchmark in its own right (Section~\ref{sec:results}), and it is the fully
controlled environment in which the mutation-audit methodology (Section~\ref{sec:audit-method}) was
built and certified before being pointed at external artifacts. The benchmark is structured into
two distinct evaluation tracks:

\begin{itemize}[leftmargin=*]
  \item \textbf{Track A (Static Generation):} Tasks a model with generating a synthesizable implementation de novo from a natural-language specification and a rigidly defined port interface.
  \item \textbf{Track B (Agentic Repair):} Tasks an agent with optimizing a pre-existing, functionally correct baseline repository to achieve a specific engineering objective (e.g., achieving timing closure, reducing area, eliminating an inferred latch, or implementing clock-domain-crossing safety for a FIFO). The agent operates under a strict token, wall-clock, and tool-call budget, and is strictly prohibited from modifying the testbench, formal properties, timing constraints, or task metadata.
\end{itemize}

Both tracks utilize a unified scoring pipeline consisting of a deterministic sequence of verification gates followed by a Power, Performance, and Area (PPA) computation. Stages 0--2 enforce pass/fail functional criteria: Verilator linting, Icarus/cocotb simulation against held-out hidden test vectors, and bounded model checking via SymbiYosys (for tasks defining formal properties). Crucially, functional correctness is enforced as a strict prerequisite rather than a weighted parameter (ADR-0000). A design failing any component of Stages 0--2 receives a score of zero, irrespective of its physical metrics. This mechanism eliminates the most prevalent failure mode of PPA-weighted scoring: adversarial optimization, wherein a model deliberately truncates functionality to minimize synthesized area. Only designs that successfully clear every correctness gate advance to Stages 3--5, which comprise Yosys synthesis targeting the sky130hd standard-cell library, OpenSTA static timing analysis at the specified task clock target, and power estimation. Track B imposes an additional verification gate utilizing sequential equivalence checking (\texttt{eqy}) against the baseline design for objectives declared behavior-preserving (Section~\ref{sec:trackb} details which objectives this covers)---ensuring an agent cannot achieve timing or power closure via covert behavioral modification---and mandates immediate disqualification if the agent modifies the test harness, formal properties, timing constraints, or metadata (files orthogonal to legitimate PPA or timing optimization).

For designs that successfully clear all verification gates, the final score is derived as a capped geometric mean relative to the human-reviewed reference implementation:
\begin{verbatim}
ppa        = geomean(ref_area / area, ref_delay / delay, ref_power / power)
task_score = 100 * min(ppa, 1.5) / 1.5
\end{verbatim}
In this formulation, $delay$ denotes the worst-path delay at the task's clock target as reported by OpenSTA. The reference metrics ($ref\_area$, $ref\_delay$, $ref\_power$) are extracted from the human-reviewed implementation, which intrinsically yields a base PPA score of 1.0. The 1.5 cap is imposed to mitigate degenerate outlier exploitation: without it, an extreme improvement concentrated in a single axis (e.g., an area ratio orders of magnitude above baseline, with delay and power essentially unchanged) can dominate the geometric mean and yield an unbounded score that overstates a design whose other two dimensions never meaningfully improved. This cap establishes a practical optimization ceiling of 1.5$\times$ relative to the human baseline, beyond which further optimization does not yield additional score improvements. The cap bounds the score from above; it does not by itself prevent a poor result on one axis from being averaged against strong results on the other two, which is an inherent property of any geometric-mean aggregation and is not specific to this cap. The Track A leaderboard reports Pass@1, mean PPA (calculated exclusively over passed tasks), and a composite GateTruth Score (Pass@1 multiplied by normalized mean PPA), stratified by complexity tier. Track B reports the objective-met rate, median PPA delta, and median operational cost (measured in tokens, wall-clock seconds, and tool calls). Cost is presented as a primary reporting metric, reflecting the principle that an agent achieving an objective via unbounded resource expenditure fails to demonstrate practical utility.

The task suite encompasses three complexity tiers in addition to the Track B agentic dataset: Tier 1 (20 tasks) covers combinational logic and fundamental sequential building blocks (e.g., priority encoders, Gray-code converters, barrel shifters, one-hot FSMs); Tier 2 (25 tasks) involves protocol and datapath components (e.g., UART, SPI, AXI4-Lite, round-robin arbiters, FIFOs); and Tier 3 (15 tasks) focuses on pipelined arithmetic, localized accelerators, and memory-system fragments. Track B comprises 8 agentic objectives layered over existing task repositories. Each task is distributed as a standardized, immutable package: an original-prose specification, a locked interface, a human-reviewed reference implementation, a bifurcated testbench (public smoke test and hidden scoring vectors), an applicable formal property subset, a singular timing constraint file, and machine-readable metadata. This uniform schema guarantees that novel task authoring and submission grading operate through identically defined pipeline infrastructure.

\section{Task Suite}
\label{sec:suite}

The GateTruth Track A suite comprises 60 specification-to-RTL tasks categorized into three tiers based on functional and physical complexity, intentionally selected to represent a realistic distribution of digital design challenges rather than clustering at a single difficulty level. Tier 1 (20 tasks) encompasses combinational logic and fundamental sequential building blocks, including priority encoders, binary-to-one-hot and Gray-code converters, barrel shifters, leading-zero counters, one-hot finite state machines, and parameterizable counters. Tier 2 (25 tasks) involves protocol and datapath components at the intersection of control and data flow, such as UART transmit and receive modules, SPI master and slave interfaces, an AXI4-Lite register file, round-robin arbiters, synchronous FIFOs, clock-domain-crossing synchronizers, and stream width converters. Tier 3 (15 tasks) focuses on pipelined arithmetic, localized accelerators, and memory-system fragments, including a Booth multiplier, a pipelined multiplier, a loadable FIR filter, a first-order IIR filter with genuine feedback, a sequential divider, a systolic processing-element tile, a cache-tag comparator, and a CRC-32 datapath. Additionally, 8 Track B agentic tasks are derived from existing task repositories; each couples an intentionally suboptimal yet functionally correct baseline design with a specific engineering objective (Section~\ref{sec:trackb}).

The primary utility of the suite derives from a uniform authoring standard that ensures every task constitutes a fair, reproducible, and contamination-resistant measurement. Each task is structured as a rigidly defined seven-part package: (1) an original-prose specification embedding a unique canary GUID; (2) a locked port interface that establishes a fixed module boundary, ensuring submissions are directly comparable and synthesizable against a uniform testbench; (3) a human-reviewed reference implementation---maintained as a draft until formal maintainer sign-off---which intrinsically defines the PPA denominator with a normalized score of exactly 1.0; (4) a bifurcated testbench comprising a public smoke-test section and a designated hidden scoring section; (5) a formal-property subset applicable to tasks where functional correctness is amenable to bounded model checking (46 of the 60 Track A tasks, 77\%); (6) a singular timing-constraint file specifying a dedicated clock target for the task; and (7) machine-readable metadata documenting the tier, descriptive tags, the clock target, formal verification applicability, and dual human-review sign-off fields. Because this package schema is utilized identically for both task authoring and submission grading, the evaluation pipeline remains tightly coupled to the authoring pipeline, preventing methodological drift as the suite expands.

Two foundational authoring principles are critical to the validity of the benchmark. First, a task-specific clock target is strictly enforced, explicitly rejecting a universal timing constraint across the suite. For instance, a 16$\times$16-into-48-bit multiply-accumulate operation inherently necessitates a more relaxed clock period than a simple barrel shifter. Imposing a uniform target would invariably render fundamental tasks trivially timed or complex tasks unsatisfiable, thereby degrading the discriminatory utility of the timing metric. Consequently, 56 tasks target a 10~ns period, while the four most computationally intensive datapaths define relaxed targets (12--20~ns), specifically sized to their critical paths and explicitly documented in the task metadata. Second, the reference implementation functions as a rigorous human sign-off gate rather than a purely automated artifact. While every reference is verified end-to-end via the automated framework---encompassing linting, simulation, formal verification (where applicable), and the complete synthesis-to-timing-and-power flow---it is not designated as the golden reference until a maintainer manually validates the RTL against the specification and formally authorizes it. This principle acknowledges that the reference serves as the absolute baseline against which the entire leaderboard is evaluated; a benchmark predicated on ground truth that has never undergone human review asserts a correctness claim it fundamentally cannot substantiate.

\section{Rigor}
\label{sec:rigor}

Three foundational properties of GateTruth's architecture warrant explicit articulation, as each is deliberately engineered to mitigate specific methodological vulnerabilities that can inadvertently inflate benchmark results without malicious intent. The first of these, mutation-validated testbenches, is the mechanism generalized into the external audit methodology of Section~\ref{sec:audit-method}.

\paragraph{Mutation-Validated Testbenches.} The absence of failures in a testbench is insufficient evidence of a functionally correct design; it may merely indicate a failure to stimulate critical logic paths. Consequently, every task's testbench is required to eliminate (kill) a minimum of 95\% of semantic mutants generated from the reference implementation. These mutations---comparator-boundary flips, operator inversions, logic and bitwise inversions, shift-direction inversions, reset- and enable-polarity flips, assignment deletion and register-hold mutations, and output inversions---are applied deterministically under a fixed seed, using a structural generator that masks comments and string literals prior to matching and enumerates every viable code site (a naive literal-substitution approach previously left many tasks generating zero mutants, a vacuously passing gate). During M1 development, three pilot tasks were validated using this protocol, each achieving a 100\% kill rate (gray-code counter: 8/8 mutants; synchronous FIFO: 13/13; UART transmitter: 14/14). Crucially, every recorded kill was attributed to functional simulation failures rather than linting errors; each mutant compiled cleanly and was rejected due to observed incorrect behavior, rather than syntactic breakage. This distinction is critical, as a mutation gate that artificially reaches 100\% by rejecting non-compiling mutants at the lint stage provides zero signal regarding testbench quality. To adversarially validate the stringency of the gate itself, a synthetic task featuring a deliberately vacuous, always-passing testbench was evaluated; the mutation harness correctly registered a 0\% kill rate and failed the gate, confirming that the $\ge$95\% threshold is rigorously enforceable rather than trivially satisfiable.

The gate now encompasses the entire evaluation suite. Under deterministic, single-threaded, fixed-seed execution, all 60 Track A testbenches kill at least 95\% of their generated mutants: 56 achieve 100\%, while the remaining four (an AXI4-Lite register file at 97.6\%, 41/42 mutants killed; an I2C slave and a UART receiver at 97.1\%, 34/35 each; and a Booth multiplier at 96.8\%, 30/31) each retain a single genuine survivor above the 95\% floor, properly documented in the suite's mutation log. One methodological observation from this certification process warrants reporting, as it represents a vulnerability applicable to any mutation-gated benchmark, external or our own: mutation verdicts are not invariant to execution concurrency. A mutant whose simulation straddles the per-run time budget may be classified as ``killed by timeout'' under heavy parallel load, yet ``survived'' when executed sequentially, allowing parallel sweeps to silently inflate kill rates. Therefore, GateTruth certifies verdicts exclusively from sequential runs and classifies simulation timeouts as indeterminate---penalizing the kill rate rather than incrementing it---with a single extended-budget retry utilized for disambiguation. This sequential-execution discipline is enforced identically whether the mutation engine is measuring our own suite (this section) or an external benchmark (Section~\ref{sec:audit-method}).

\paragraph{Contamination Defense.} GateTruth implements a multi-layered defensive architecture rather than relying on a singular mechanism, operating on the premise that while individual defenses can be circumvented, their aggregation raises the cost of adversarial exploitation beyond the cost of legitimately solving the task. Public testbenches contain only a smoke-test section; the scoring vectors determining definitive pass/fail status are isolated under a marked \texttt{\#\ -{}-{}- HIDDEN -{}-{}-} section during development and are extracted to a permanently private repository at the v1.0 freeze. Thus, a model that memorizes a public specification-testbench pair remains blind to the evaluation vectors. Furthermore, every specification utilizes original prose---expressly prohibiting incorporation from HDLBits, VerilogEval, textbooks, or public repositories---and embeds a task-unique canary GUID. The canary's evidentiary value is prospective, not immediate: a model asked to solve a given task is trivially given that task's own canary in its prompt, so reproducing it there is not itself evidence of anything. The canary instead lets us detect \emph{future} contamination --- if a public specification is later scraped into a model's training corpus, that model's outputs on \emph{unrelated} prompts, or its behavior when queried about the canary string outside the context of its own task, becomes attributable to this specific released spec rather than to generic RTL knowledge. For v1.0, submission integrity rests on maintainer-executed scoring: every leaderboard entry is produced by the maintainer running the candidate model through the pinned reference image, so no externally-reported result is trusted. A versioned, pull-request-based submission model, in which contributors self-report results and the maintainer independently replicates a random sample against the pinned image, is documented as a roadmap item rather than a v1.0 capability. A more robust architecture---incorporating an asymmetric public/hidden task split with published hash commitments and semantic-preserving inference-time specification mutation---is likewise deliberately deferred to v1.1 (ADR-0002) rather than deployed incompletely for the initial release.

\paragraph{Pinned-Digest Reproducibility.} Every evaluated run executes strictly within a singular Docker container image. The resulting result manifest records the container's content digest alongside exhaustive stage-by-stage pipeline telemetry, including lint warnings, simulation pass counts, formal verification status, synthesized area and cell count, worst-negative-slack, achieved frequency, and estimated power. Successive evaluations of identical submissions against the same task and image digest yield byte-identical canonical JSON output (excluding timestamp, signature, and wall-clock fields). This deterministic behavior was verified end-to-end during the M1 phase (flow-determinism ticket, SB-007) and holds universally for all subsequently authored tasks. The manifest schema enforces this determinism as a rigid contract, not a best-effort property: any evaluation flow yielding nondeterministic synthesis or timing output is classified as a defect within the flow scripts, not an accepted variance in leaderboard noise. Consequently, given the pinned image digest, the submission file, and (for correctness-gate stages) the same hidden test vectors, reproducing any published score depends on no other input: it requires no wall-clock-sensitive tool behavior and no manual intervention. This is a claim about the pipeline's own determinism, not about third-party access to the hidden vectors themselves, which remain privately held for the contamination-resistance reasons described above; a party without those vectors can independently reproduce the deterministic, non-scoring pipeline stages (lint, synthesis, timing, power) but not the private correctness-gate verdict.

\section{Auditing External Benchmarks}
\label{sec:audit-method}

The mutation-certification protocol of Section~\ref{sec:rigor} generalizes directly: nothing about
injecting deterministic, seeded mutants into a reference implementation and measuring what fraction
a testbench kills is specific to GateTruth's own tasks. This section describes the audit
methodology applied, unmodified in its core mutation-generation and verdict logic, to external RTL
benchmarks; Section~\ref{sec:audit-results} reports the resulting measurements.

\paragraph{Scope and target selection.} We selected RTLLM v2.0~\cite{lu2024rtllm,liu2025openllmrtl}, a 50-design
RTL-generation benchmark whose public repository is MIT-licensed and whose testbenches are
self-checking Verilog runnable under an open-source simulator, and NVIDIA's CVDP~\cite{pinckney2025cvdp},
whose harness is Apache-2.0 and whose public dataset mixes licenses by content type: non-code
material under CC~BY~4.0 and code under Apache-2.0, per the dataset's own NOTICE file. Both harness
and dataset licenses permit the read-only measurement this audit performs. We deliberately excluded
any benchmark whose license terms were ambiguous or unverified at audit time.

\paragraph{Read-only, pinned provenance.} Vendor repositories and dataset revisions are fetched
once, pinned to an exact commit hash or dataset revision, and treated as read-only for the
remainder of the audit; our own outputs are written only to a separate results directory and never
back into the vendor tree. Every reported measurement records the exact pinned commit or revision
it was produced against, so the audit is independently reproducible by a third party against the
same vendor snapshot.

\paragraph{Mandatory baseline validation.} Before any mutant is generated for a given design, the
\emph{unmodified} reference implementation must pass the benchmark's own testbench under our
harness. A design that does not clear this bar --- because its testbench requires a commercial
simulator we do not run, or a harness convention our adapter does not yet support --- is reported
with an explicit \texttt{unsupported} status and excluded from kill-rate aggregation. It is never
scored as a misleading 0\% kill rate, which would conflate ``the testbench caught nothing'' with
``we could not run the testbench at all.''

\paragraph{Build-convention normalization, not silent inflation.} Some external benchmarks ship a
golden reference whose internal module name differs from the name their own testbench instantiates
(for instance, a file that declares one module identifier while the benchmark's own per-design build
recipe expects the candidate file to be renamed and re-moduled before compilation). Naively compiling
such a design against its testbench fails not because the design is genuinely incompatible with an
open-source simulator, but because our harness has not yet replicated the vendor's own build
convention. Where a design's own build recipe documents this convention, the audit reproduces it at
compile time only: the golden source is copied into a temporary location, its module declaration is
rewritten to the name the testbench expects, and only that temporary copy is compiled --- the pinned
vendor file on disk is never modified. Every report row records whether this alias was applied and,
if so, the exact rename performed, so a third party can distinguish ``passed as shipped'' from
``passed after a documented, vendor-convention-matching alias'' at a glance. Designs that still fail
to compile or run after this normalization are reported as genuine incompatibilities, not silently
retried with looser semantics.

\paragraph{Operator set.} The audit uses exclusively the generic, benchmark-agnostic mutation
operators already certified against GateTruth's own suite in Section~\ref{sec:rigor} (comparator
and operator inversions, logic and bitwise inversions, shift-direction inversion, reset- and
enable-polarity flips, assignment deletion and hold, output inversion). It does not use any
mutation specification authored for a specific GateTruth task, since those are not generic and
would not be a fair test of an external testbench's own rigor.

\paragraph{Execution policy.} Identical to Section~\ref{sec:rigor}: a fixed seed determines mutant
ordering; execution is strictly sequential (never parallel) for any measurement reported as final;
a mutant that times out is retried once at an extended budget, and a mutant that still times out is
recorded as indeterminate and counted against the kill rate, never simply dropped. The full sweep is
run once at a fixed seed; to certify that this is not a one-off artifact of that run, a
deterministically selected sample of the eligible designs (Section~\ref{sec:audit-results} gives
the exact sampling rule and the selected designs) is independently re-run at the same seed, and
every sampled design's per-design result must be byte-identical across the two runs before the full
sweep is treated as certified. This is a sampling-based determinism check, not a claim that every
individual design was executed twice.

\paragraph{The CVDP gap, reported rather than hidden.} NVIDIA's public CVDP release deliberately
withholds reference solutions (the \texttt{output}/\texttt{patch} fields) specifically to mitigate
data contamination in model evaluation. This is a reasonable design choice for CVDP's own purpose
(scoring generated code against a held-out answer), but it removes exactly the artifact --- a known
correct golden design --- that mutation testing requires as its starting point. Where our own
feasibility sweep confirms no usable golden RTL exists in the public release, we report that as a
structural finding about the benchmark's auditability (Section~\ref{sec:audit-results}), rather than
silently omitting CVDP from this paper or fabricating a workaround.

\section{External Audit Findings}
\label{sec:audit-results}

This section reports mutation-kill measurements for RTLLM v2.0's public testbenches, produced by
the methodology of Section~\ref{sec:audit-method} against vendor commit
\texttt{41b26896e33b536940116a975626455eed3de65e} fetched on 2026-07-29, compiled and simulated
with Icarus Verilog 12.0 (stable, \texttt{v12\_0}). Of the 50 designs in RTLLM v2.0, 44 passed
baseline validation and were audited (26 required the compile-time module alias of
Section~\ref{sec:audit-method}, 18 passed as shipped); 6 were reported \texttt{unsupported}. Four of the six are genuine
Icarus Verilog-2001 incompatibilities in the vendor's own shipped files: \texttt{asyn\_fifo}'s
testbench uses a SystemVerilog \texttt{break} statement, \texttt{clkgenerator} and
\texttt{freq\_divbyodd} each drive a \texttt{reg} from a continuous assignment in a way Icarus
rejects, and \texttt{multi\_8bit}'s reference file contains a \texttt{for}-loop syntax error. The
remaining two are more consequential findings in their own right rather than tooling gaps:
\texttt{ring\_counter}'s testbench requires a SystemVerilog whole-array assignment, and, most
notably, \texttt{radix2\_div}'s \emph{unmodified golden reference fails its own shipped
testbench} (three of its self-checking vectors report a mismatch), independent of anything
GateTruth's mutation engine does.

Across the 44 audited designs, the median mutation-kill rate was 71.2\%, and 32 of the 44---73\%---
fall below the 95\% floor every one of GateTruth's own 60 Track A tasks clears (Section~\ref{sec:rigor}).
Three designs (\texttt{adder\_8bit}, \texttt{edge\_detect}, \texttt{square\_wave}) score a 0\%
kill rate: none of the generic mutants injected into their reference implementations were caught
by RTLLM's own shipped testbench. Their mutant counts are small (1, 10, and 6 respectively, versus a median of 10.5 mutants per
design across all 44, ranging from 1 to 115) --- \texttt{adder\_8bit} in particular
generates only a single viable generic mutant from its reference source, so its 0\% figure should
be read as a single-mutant miss, not as evidence the testbench misses a broad class of faults; the
other two 0\% results rest on a somewhat larger but still modest number of mutants. Full per-design
mutant counts, operator breakdowns, and pass/fail logs for all 44 audited and 6 unsupported designs
are committed at \texttt{external-audit/results/rtllm/final/*.json} and summarized in
\texttt{external-audit/results/summary.md} in the GateTruth repository
(\url{https://github.com/meetbhadra701-cloud/GateTruth}, currently private during review and made
public upon submission), so every number in this section can be checked against the raw evidence.
At the other extreme, 12 of the 44 designs do
reach a 100\% kill rate, so the result is not that RTLLM's testbenches are uniformly weak, but that
their rigor is highly uneven and, for a majority of designs, well short of what a mutation-kill
floor would require. A deterministic re-run of a fixed one-in-five sample (8 of the 44 eligible
designs, 18.2\%, selected via \texttt{random.Random(20260729).sample}) reproduced byte-identical
per-design results across two independent runs, certifying the measurement.

For CVDP, the Day-1 feasibility sweep (Section~\ref{sec:audit-method}) found 0 usable golden
designs out of the 302 rows inspected in the public non-agentic, no-commercial-tooling subset:
every row's \texttt{output} field is withheld as shipped. Seventeen rows carry an OSS-origin
provenance tag in the dataset's NOTICE file, but the pinned public harness supplies no verified,
overlapping golden RTL for those rows, so they are not counted as usable and no CVDP mutation-kill
measurement is reported. This is treated as a structural finding about the current public
release's auditability, not a shortfall of this paper's methodology: a mutation-kill audit
requires a known-correct starting point to mutate, and the public CVDP release does not supply
one, by NVIDIA's own design, for anti-contamination reasons.

\section{Track B: Agentic Protocol}
\label{sec:trackb}

Track B evaluates a fundamentally distinct capability compared to Track A: rather than assessing de novo generation from a specification, it measures the capacity of an autonomous agent to optimize an existing, functionally correct design toward a concrete engineering objective. These objectives include achieving timing closure at a target frequency, reducing area by a specified percentage, eliminating inferred latches, or ensuring clock-domain-crossing (CDC) safety for a FIFO---all while preserving orthogonal system functionality. This paradigm more closely approximates the operational realities of RTL codebase maintenance compared to isolated generation tasks. Furthermore, it exposes a distinct adversarial failure mode: under strict budget constraints, an agent is incentivized to artificially satisfy the literal parameters of an objective by modifying the evaluation testbench rather than the device under test (DUT).

To mitigate this, the evaluation protocol is strictly constrained. An agent is provisioned with a checked-out repository alongside an \texttt{objective.yaml} file that explicitly defines the engineering target, an explicit \texttt{behavior\_preserving} flag, and a strict budget measured in tokens, wall-clock seconds, and tool calls. Within the execution container, the agent is restricted to standard CLI commands (\texttt{sb-lint}, \texttt{sb-sim}, \texttt{sb-synth}, and \texttt{sb-sta}), and its sole permitted output artifact is the modified repository. The evaluation harness enforces an unconditional, rigidly ordered post-processing sequence: the complete test suite must execute successfully; where the objective is declared behavior-preserving, sequential equivalence checking (\texttt{eqy}) against the pre-modification baseline design must additionally pass; and only upon clearing these prerequisites is the stated objective evaluated and the PPA delta and computational cost recorded. Three of the eight tasks (the two timing-closure objectives and the power-reduction objective) are declared behavior-preserving and gated by \texttt{eqy}. Four of the remaining five --- latch removal, CDC safety, arbiter fairness, and AXI byte-enable handling --- are declared \texttt{add\_property} objectives whose entire point is to change the design's observable behavior on previously out-of-specification inputs, so they are checked against the hidden testbench (which is updated to cover the new property) rather than against sequential equivalence with the unmodified baseline. The area-reduction objective is the fifth: despite area reduction ordinarily being a behavior-preserving optimization, this task is explicitly flagged \texttt{behavior\_preserving: false} in its own task definition and is therefore not \texttt{eqy}-gated. It is not exempt from correctness checking --- Stage~1's hidden-testbench simulation, using the same mutation-certified vectors as the underlying Track A task, still runs unconditionally for every Track B submission regardless of this flag, so a truncate-the-logic attack would still have to survive that same test vector set --- but it does not receive \texttt{eqy}'s stronger, exhaustive equivalence guarantee. We treat this as a genuine, disclosed gap in this task's design rather than paper it over: extending \texttt{eqy} coverage to the area- and power-reduction objectives is future work (Section~\ref{sec:limitations}).

The sequential-equivalence gate serves as the foundational mechanism for the behavior-preserving objectives. It formally proves that the agent's modified design implements the identical logical function as the original baseline, thereby precluding adversarial PPA or timing manipulation (e.g., logic truncation to minimize synthesis area or critical path delays) on tasks where behavior is not supposed to change at all. Additionally, any modification to the testbench, formal properties, timing constraints, or task metadata results in immediate task disqualification, which is explicitly flagged in the result manifest, regardless of which gate applies to the objective. Because a legitimate structural optimization or property addition should never necessitate modifications to these specific validation files, such diffs are interpreted as a definitive signal of an invalid agent approach rather than a borderline case requiring manual adjudication.

Finally, an agent that exceeds its allocated budget is forcibly terminated and evaluated strictly on its state at termination; consequently, unbounded-cost ``successes'' are not rewarded via omission from the leaderboard's cost reporting. Cost metrics---specifically median tokens, wall-clock seconds, and tool calls---are reported explicitly as independent variables rather than being aggregated into a single opaque score. This reporting structure preserves the critical trade-off matrix between objective satisfaction and the computational expenditure required to achieve it, providing practitioners with the exact empirical data necessary to evaluate the viability of deploying such agentic workflows.

\section{Reference-Suite Baseline Results}
\label{sec:results}

Having established GateTruth's own suite as the certified environment in which the audit
methodology was built (Sections~\ref{sec:rigor}--\ref{sec:audit-method}), we report its baseline
results as a secondary, independently useful measurement: how do current models and agents
actually perform against this suite? We present the empirical results of the v1.0 baseline evaluation campaign, run 2026-07-22, encompassing seven models evaluated in that campaign (not necessarily each provider's current frontier at time of reading): three from Anthropic (Claude Opus 4.8, Sonnet 4.6, Haiku 4.5), two from OpenAI (GPT-5, GPT-5-mini), and two accessed via OpenRouter (Google Gemini 2.5 Pro, Meta Llama 4 Maverick). These models were evaluated across the complete suite of 60 Track A tasks and 8 Track B tasks. All evaluations utilized a single-sample inference protocol at a temperature of zero, substituting provider-default sampling exclusively for reasoning models that do not support explicit temperature parameterization. Executions were strictly confined within a singular, content-addressed container environment (\texttt{gatetruth:v1}, digest recorded per manifest) under a \$300 per-run spend cap, and scored exclusively against maintainer-reviewed reference RTL and held-out hidden vectors. Accounting for both tracks across all seven models, and reconciled against the settled provider billing ledger (Anthropic \$18.79, OpenAI \$14.54, OpenRouter \$17.75), the campaign's total metered cost was \$51.08; this reconciled total includes transient-transport retries and is not directly derivable by summing the per-model costs in Tables~\ref{tab:tracka}--\ref{tab:trackb}, which report only the cost of each model's accepted official sample. This \$51.08 figure covers only the 2026-07-22 v1.0 campaign itself; the Appendix~\ref{app:variance} run-to-run variance study is a separate, later evaluation (2026-07-26) that cost an additional \$8.18 and is not included in this total. Given the recorded container digest and a specific submission's RTL text, re-scoring that submission is byte-level reproducible; this reproducibility claim applies to the scoring pipeline, not to re-invoking the models themselves, whose sampling is not deterministic (Appendix~\ref{app:variance} quantifies the resulting run-to-run variance).

For Track A (Table~\ref{tab:tracka}), the models are ranked by GateTruth Score, defined in Section~\ref{sec:design} as Pass@1 multiplied by normalized mean PPA; equivalently, because failed tasks contribute a task score of exactly zero, this is the mean capped-PPA task score across all 60 tasks on a 0--100 scale, where the human-authored reference implementation is normalized to $100/1.5 = 66.\overline{6}$ (displayed as 66.7). The resulting hierarchy is: Opus (46.91), Sonnet (45.03), Haiku (33.58), GPT-5-mini (33.44), GPT-5 (31.64), Llama 4 Maverick (25.52), and Gemini 2.5 Pro (17.83). These are single-run figures, and Opus and Sonnet sit within run-to-run variance of one another at the top of the table (Appendix~\ref{app:variance}); a clear gap then separates this frontier pair from the remaining models. The correctness-gate pass@1 metric exhibits a corresponding monotonic decline: Opus 42/60, Sonnet 41/60, Haiku 31/60, GPT-5-mini 30/60, GPT-5 28/60, Llama 4 Maverick 23/60, and Gemini 2.5 Pro 16/60. A critical observation from this data is the substantial performance deficit at the frontier: the highest-performing model (46.91) achieves approximately 70\% of the human reference baseline (66.7). This indicates that even state-of-the-art systems forfeit a measurable margin of physical quality compared to a competent human implementation. Furthermore, two structural insights emerge regarding the evaluation paradigm. First, inference cost and physical quality are decoupled in specific, checkable ways rather than uniformly: Llama 4 Maverick's score (25.52) is 43\% above the field's lowest score while its metered cost (\$0.029) is roughly 7$\times$ below its next-cheapest competitor (GPT-5-mini, \$0.203) and well over an order of magnitude below the three most expensive models (Opus, GPT-5, Gemini); Gemini 2.5 Pro is simultaneously the most expensive model evaluated and the lowest-scoring; and GPT-5 underperforms relative to its smaller counterpart, GPT-5-mini. Such inversions remain entirely opaque to conventional pass@k benchmarks. Second, on Track A the aggregate ranking is governed primarily by the correctness gate---how many tasks each model passes---with post-synthesis PPA contributing a small, signed adjustment rather than a reordering: measured only over the tasks a model passes, the capped geometric mean neither reorders the table nor reliably separates the two frontier models, which are statistically indistinguishable on Track A across repeated runs (Appendix~\ref{app:variance}). Single-shot generation therefore separates current models mainly by correctness; the physical-quality axis, while genuinely measured here, becomes a decisive differentiator only in the agentic-repair track (Section~\ref{sec:trackb}), where optimizing a design against an explicit PPA objective is the task itself.

\begin{table}[t]
\centering
\caption{Track A: static specification-to-RTL generation, all 60 tasks. GateTruth Score is the mean capped-PPA task score on a 0--100 scale; the human reference is normalized to 66.7 by construction.}
\label{tab:tracka}
\small
% generated-on: 2026-08-05 git-sha: b71ea2f897a08ed1721f9a8ff5f6fde45710c58c
\begin{tabular}{llrrrr}
\toprule
Provider & Model & Tasks & Score & Tokens & Cost (USD) \\
\midrule
anthropic & claude-opus-4-8 & 60 & 46.91 & 185879 & 1.578055 \\
anthropic & claude-sonnet-4-6 & 60 & 45.03 & 149298 & 0.855978 \\
anthropic & claude-haiku-4-5-20251001 & 60 & 33.58 & 151785 & 0.298001 \\
openai & gpt-5-mini & 60 & 33.44 & 171464 & 0.202953 \\
openai & gpt-5 & 60 & 31.64 & 145635 & 0.925733 \\
openrouter & meta-llama/llama-4-maverick & 60 & 25.52 & 122387 & 0.028870 \\
openrouter & google/gemini-2.5-pro & 60 & 17.83 & 303095 & 2.164052 \\
\bottomrule
\end{tabular}

\end{table}

Track B (Table~\ref{tab:trackb}) serves as a substantially more discriminating evaluation instrument. Measured by the objective-met rate across the 8 agentic-repair tasks, performance degrades precipitously: Opus achieved 5/8, Sonnet 3/8, GPT-5 1/8, and Haiku, GPT-5-mini, Llama 4 Maverick, and Gemini 2.5 Pro failed to satisfy any objectives (0/8). Consequently, the imposition of agentic PPA optimization under a strict sequential-equivalence constraint isolates model capabilities far more aggressively than de novo generation: of the seven evaluated models, one (Opus) cleared a majority of the repair objectives, one (Sonnet) cleared more than a third, one (GPT-5) cleared a single objective, and four failed entirely. This pattern is consistent with the concern raised in Section~\ref{sec:intro} that syntactically valid generation and genuine engineering repair are distinct capabilities; we do not treat this as a formally pre-registered hypothesis. No model successfully satisfied either of the two timing-closure objectives (optimizing a tightened clock constraint on a multiply-accumulate unit and a multiplier) or the IIR power-reduction objective --- that is, all three behavior-preserving, \texttt{eqy}-gated objectives went unsolved by every model. Successful interventions were confined to the five remaining objectives: removing an inferred latch, implementing clock-domain-crossing safety for a FIFO, correcting arbiter fairness, rectifying AXI byte-enable handling, and reducing the FIR filter's area by at least 30\%. Opus is the only model to clear all five of these (its 5/8 total); Sonnet clears three of the four \texttt{add\_property} objectives (latch, CDC, arbiter, but not AXI byte-enable); and GPT-5's sole completion is the AXI byte-enable task. These findings suggest that timing closure under constrained parameters---a repair objective necessitating post-synthesis critical path reasoning rather than the syntactic pattern completion of established templates---remains unsolved among the models evaluated in this campaign (Section~\ref{sec:results}).

\begin{table}[t]
\centering
\caption{Track B: agentic PPA repair, 8 tasks. Objective-met counts require the Stage-1 correctness gate, plus \texttt{eqy} sequential-equivalence preservation for the 3 behavior-preserving objectives (Section~\ref{sec:trackb}). Area and power ratios are candidate/baseline (lower is an improvement). WNS delta is candidate WNS minus baseline WNS (a positive delta is more slack, i.e., an improvement). All three PPA-delta figures are the per-model median taken only over the 4 objectives that report a quantitative PPA delta at all (the two timing-closure tasks, area reduction, power reduction) and that reached the correctness gate; the other 4 \texttt{add\_property} objectives (latch removal, CDC safety, arbiter fairness, AXI byte-enable) never report a PPA delta, pass or fail, since their objective is not a PPA change. A value of exactly $1.0000\text{x}$/$1.0000\text{x}$/$+0.000$~ns means the model's candidate for at least one of those 4 tasks reached the correctness gate but produced no net physical change; \texttt{N/A} means none of the 4 ever reached the correctness gate for that model.}
\label{tab:trackb}
\small
\setlength{\tabcolsep}{4pt}
\resizebox{\textwidth}{!}{% generated-on: 2026-07-24 git-sha: abca3d57f8ed
\begin{tabular}{llrrlr}
\toprule
Provider & Model & Objectives met & Objective-met rate & Median PPA delta & Cost (USD) \\
\midrule
anthropic & claude-opus-4-8 & 5/8 & 62.50\% & area=0.5527x; power=0.3687x; WNS=-1.555 ns & 1.647805 \\
anthropic & claude-sonnet-4-6 & 3/8 & 37.50\% & area=1.0000x; power=1.0000x; WNS=+0.000 ns & 4.157709 \\
openai & gpt-5 & 1/8 & 12.50\% & area=1.0000x; power=1.0000x; WNS=+0.000 ns & 3.955635 \\
anthropic & claude-haiku-4-5-20251001 & 0/8 & 0.00\% & area=1.0000x; power=1.0000x; WNS=+0.000 ns & 2.043547 \\
openrouter & google/gemini-2.5-pro & 0/8 & 0.00\% & area=1.0000x; power=1.0000x; WNS=+0.000 ns & 6.593550 \\
openai & gpt-5-mini & 0/8 & 0.00\% & area=N/A; power=N/A; WNS=N/A & 0.503571 \\
openrouter & meta-llama/llama-4-maverick & 0/8 & 0.00\% & area=N/A; power=N/A; WNS=N/A & 0.041961 \\
\bottomrule
\end{tabular}
}
\end{table}

Synthesizing the data across both tracks reveals a consistent ordinal ranking, though Track B significantly amplifies the performance delta between models. Systems capable of generating physically superior silicon are concomitantly more proficient at processing synthesis-and-timing feedback to execute targeted repairs. The performance gap between the top-tier models and the remainder of the field is notably wider in agentic repair than in initial generation. This divergence validates a secondary objective of the reference suite: to differentiate between models that merely emit syntactically simulatable RTL---a baseline capability shared by most evaluated models---and those capable of producing, or iteratively refining, silicon to a standard acceptable to a competent hardware engineer.

\section{Limitations and Threats to Validity}
\label{sec:limitations}

\paragraph{Audit scope.} The external audit (Sections~\ref{sec:audit-method}--\ref{sec:audit-results}) covers RTLLM v2.0 in full and reports a structural finding, rather than a measurement, for CVDP given its withheld reference solutions. It does not cover every RTL-generation benchmark in current use. In particular, VerilogEval~\cite{liu2023verilogeval,pinckney2024revisiting} --- arguably the field's most-cited benchmark --- is not audited here; this is a scope decision made to keep the first application of this methodology focused on one primary target (RTLLM v2.0) and one structural finding (CVDP) rather than a principled exclusion, and we name it explicitly rather than let the omission pass silently. ChipBench, ArchXBench, SLDB, and ProtocolLLM are likewise out of scope for this round. Coverage of the Si2 LLM Benchmarking Coalition's evolving CVDP-derived leaderboard is necessarily a snapshot rather than a standing measurement. The audit measures testbench sensitivity to a fixed, generic operator set; a testbench could in principle be robust to our specific operators while remaining weak against fault classes we do not inject. The methodology and tooling are open specifically so others can extend the operator set or apply it to additional benchmarks. The build-convention normalization described in Section~\ref{sec:audit-method} is bounded strictly to what a target's own documented build recipe already specifies (e.g., a per-design makefile's stated candidate-module contract); it is not a search over renamings intended to maximize the reported compatibility count, and every applied alias is recorded per design so a third party can audit the boundary between ``ran as vendor-intended'' and ``ran after our own compatibility work'' for themselves. More broadly, ``rigor'' as measured here is specifically testbench sensitivity to injected mutants; it is not a claim about a benchmark's specification correctness, golden-reference correctness, task representativeness, contamination resistance, or scoring methodology, none of which this audit examines.

\paragraph{Sensitivity is not the same as correctness.} A high kill rate establishes that a testbench is sensitive to the specific faults injected; it does not establish that the testbench's own oracle (the golden reference and the pass/fail check itself) is correct, nor that a testbench with a high kill rate would correctly accept every valid alternative implementation (specificity, or false-rejection rate, in the terminology of software mutation testing). RTLLM's \texttt{radix2\_div} illustrates the first point directly: its testbench is exercised by our mutants in the same way as any other design in the audit, but the underlying golden reference itself fails that same testbench, independent of mutation. We do not attempt to measure false-rejection rate in this work, since it would require independently correct alternative implementations or a formal specification per design, neither of which is available for RTLLM's designs; this is future work, not a property this audit currently reports.

\paragraph{Mutant generation, equivalence, and the role of the seed.} GateTruth's generator deterministically enumerates every syntactically viable mutation site for its fixed generic operator set; no sampling of mutants occurs, so the fixed seed governs only execution order (relevant to how simulation timeouts interact with per-run budgets under sequential execution, Section~\ref{sec:rigor}), not which or how many mutants are generated. We do not detect or exclude equivalent mutants (mutants that are syntactically distinct but behaviorally identical to the original under the design's reachable state space); any equivalent mutants present would depress the reported kill rate relative to a hypothetical equivalent-mutant-free count, so our reported rates are, if anything, a conservative (lower) estimate of true fault-detection sensitivity. Mutant counts per RTLLM design range from 1 to 115 (median 10.5); we report every design's raw count rather than filtering by a minimum, but results built on very small counts (as low as 1) carry correspondingly less statistical weight, and we flag this explicitly rather than let a headline percentage stand without its denominator (Section~\ref{sec:audit-results}). We also do not perform a manual audit distinguishing genuine functional-mismatch kills from incidental simulator artifacts (e.g., uninitialized-signal propagation), nor a manual audit reclassifying survivors; verdicts throughout are produced by the automated compile/simulate/pass-string pipeline described in Section~\ref{sec:audit-method}. Finally, RTLLM's 44 audited designs share common benchmark authorship, testbench-writing conventions, and a single vendor's coding style; they should not be treated as fully independent samples when interpreting the aggregate 71.2\% median or the 73\% below-floor rate.

The physical metrics reported by GateTruth's own reference suite represent synthesis and static-timing estimates rather than post-place-and-route signoff values. The evaluation flow synthesizes designs utilizing the sky130hd standard-cell library and performs timing analysis via OpenSTA at a single typical-typical process corner. Because the pipeline does not execute placement, routing, or parasitic extraction, the reported area corresponds strictly to pre-layout cell area, and the reported timing reflects pre-layout slack. While a complete place-and-route trajectory was scoped and evaluated during development, it was deferred for the v1.0 release; an open-source routed-flow feasibility assessment failed to meet the requisite reliability thresholds necessary to enforce routed metrics as a scored gate (as documented in the project's decision log). Consequently, the manuscript reports pre-route estimates---explicitly labeled as such---and designates a fully routed evaluation track as future work rather than presenting it as a silent approximation. Therefore, PPA comparisons within this paper are rigorously valid within the confines of the pinned evaluation flow, but they must not be extrapolated as predictive of taped-out physical silicon.

Two additional methodological limitations bound the reference suite's claims regarding data contamination and functional correctness. First, the contamination defense operates as a multi-layered deterrent framework---comprising original-prose specifications, per-task canaries, and hidden test vectors excluded from the public repository. This architecture significantly elevates the computational and operational cost of leaderboard manipulation without theoretically eliminating it; a sophisticated actor capable of reconstructing hidden behavior exclusively from the public smoke testbench is not provably thwarted. A more robust asymmetric-split architecture, incorporating published hash commitments and inference-time specification mutation, is deferred to subsequent iterations. Second, the functional correctness gate is intrinsically bounded by the rigor of its underlying testbenches. We mitigate this limitation --- for our own suite, exactly as we do for the external benchmarks we audit --- by enforcing a $\ge$95\% mutation-kill floor validated under sequential execution (Section~\ref{sec:rigor}). However, mutation testing inherently quantifies a testbench's sensitivity to artificially injected faults, not its completeness against all conceivable real-world defects; notably, four of our own suite's tasks clear this threshold while retaining a small, documented number of survivors above the 95\% floor, rather than achieving a perfect kill rate. Finally, the reference suite is methodologically constrained: it utilizes English-prose specifications, operates within a single-clock paradigm (wherein clock-domain crossing is modeled structurally rather than via true multi-clock metastability simulation), and is strictly scoped to the previously defined task classes.

\paragraph{Applying the same denominator scrutiny to our own suite.} The mutant-count caveat raised above for RTLLM's audited designs (Section~\ref{sec:audit-results}) applies with the same force to our own 60 Track A testbenches: their mutant counts range from 4 to 263 with a median of 12, and 17 of the 60 have fewer than 10 mutants, meaning a single survivor in those tasks necessarily fails the 95\% floor while a task with 30+ mutants can retain several survivors and still pass. We hold RTLLM to a standard our own suite shares, not a stricter one, and disclose it symmetrically rather than only where it favors our own certification claims. Separately, Track B's 8 testbenches are not yet certified under the mutation-audit protocol described in Section~\ref{sec:rigor}; only the 60 Track A testbenches carry that certification as of this writing, and extending it to Track B is future work. The area-reduction Track B objective's lack of \texttt{eqy} coverage (Section~\ref{sec:trackb}) is a related, disclosed gap in the same direction. Finally, the human-review sign-off described in Section~\ref{sec:suite} is performed by this paper's sole author; it substitutes single-maintainer scrutiny for the multi-reviewer process a larger project would use, and we do not claim the independence multiple reviewers would provide.

\section{Related Work}

The prevailing paradigm for evaluating large language models on register-transfer-level (RTL) design predominantly assesses the functional correctness of generated code. VerilogEval~\cite{liu2023verilogeval} established this framework utilizing a suite of tasks derived from the HDLBits teaching platform, evaluated via simulation testbenches under a pass@k protocol; its subsequent iteration~\cite{pinckney2024revisiting} expanded to specification-to-RTL generation, introduced in-context prompting, and implemented granular failure classification. RTLLM~\cite{lu2024rtllm} adopted a complementary approach, originally curating thirty designs encompassing arithmetic and logic circuits, each formalized as a triplet comprising a natural-language specification, a testbench, and a human-authored reference; the project's public repository has since expanded to fifty designs under the same MIT license, released by the authors as RTLLM v2.0~\cite{liu2025openllmrtl}, and is the primary target of the external audit reported in Section~\ref{sec:audit-results}. RTLLM's own evaluation flow synthesizes generated RTL and reports post-synthesis PPA (worst-negative-slack, area, and power) against its reference designs, so it is not a correctness-only benchmark; what it does not report is a mutation-kill or equivalent systematic fault-coverage measurement of its own testbenches' rigor, which is exactly the property Table~\ref{tab:related} and Section~\ref{sec:audit-results} address. These foundational benchmarks codified a critical distinction: syntactic validity diverges significantly from functional correctness. A highly capable model may reliably emit compilable code while achieving functional correctness on less than 75\% of problems. Neither VerilogEval nor RTLLM reports a mutation-kill floor for its own testbenches --- the gap this paper's audit methodology is designed to measure. Mutation testing itself predates this literature by decades, originating in software engineering with DeMillo, Lipton, and Sayward's foundational proposal to inject faults into a program and measure what fraction a test suite detects~\cite{demillo1978hints}. Prior work has already carried that idea into hardware verification, outside the LLM-benchmark literature: Huang~et~al.~\cite{huang2015mutation} apply mutation analysis to qualify hardware testbench quality, Mantra~\cite{wu2023mantra} derives RTL mutation operators from real reported bugs, and AutoBench~\cite{qiu2024autobench} proposes automated HDL testbench generation and evaluation. None of these targets an existing public LLM RTL-generation benchmark's own shipped testbenches, which is this paper's narrower and more specific contribution.

Table~\ref{tab:related} summarizes recent RTL-generation and system-level LLM benchmarks against
two properties central to this paper: whether the benchmark reports a mutation-testing (or
equivalent systematic fault-injection) rigor measurement of its own testbenches, and whether it
scores generated designs on post-synthesis PPA rather than functional correctness alone. ChipBench~\cite{yu2026chipbench}
evaluates Verilog generation, debugging, and reference-model synthesis across three task families
using directed and constrained-random testing; ArchXBench~\cite{purini2025archxbench} scales
generation difficulty across six complexity tiers from logic primitives to full accelerators;
SLDB~\cite{alvanaki2025sldb} evaluates system-level integration and configuration of accelerators
into heterogeneous SoCs rather than module-level generation; and ProtocolLLM~\cite{sheth2025protocolllm}
narrows scope to SystemVerilog generation for four communication protocols (SPI, I\textsuperscript{2}C,
UART, AXI), evaluating syntax, synthesizability, and timing fidelity. None of the four reports a
mutation-testing or systematic fault-coverage measurement of its own testbenches in its published
paper or public repository, and none reports post-synthesis PPA scoring of the RTL it evaluates ---
based on the papers, repositories, and documentation we were able to access; we did not exhaustively
audit every appendix or supplementary artifact, and report this as the most accurate characterization
available rather than a certified negative claim.

\begin{table}[t]
\centering
\caption{Related RTL/hardware LLM benchmarks. ``Mutation rigor'' = reports a mutation-testing or
equivalent fault-injection measurement of its own testbenches. ``PPA'' = scores generated designs
on post-synthesis area/timing/power, not correctness alone.}
\label{tab:related}
\small
\begin{tabular}{lp{4.7cm}cc}
\toprule
Benchmark & Scope & Mutation rigor & PPA \\
\midrule
VerilogEval~\cite{liu2023verilogeval,pinckney2024revisiting} & HDLBits-derived generation, pass@k & No & No \\
RTLLM v2.0~\cite{lu2024rtllm,liu2025openllmrtl} & 50-design generation, self-checking TBs & No & Yes \\
CVDP~\cite{pinckney2025cvdp} & Generation + comprehension, agentic \& non-agentic & No & No \\
ChipBench~\cite{yu2026chipbench} & Generation, debugging, reference-model synthesis & No & No \\
ArchXBench~\cite{purini2025archxbench} & Tiered-complexity generation & No & No \\
SLDB~\cite{alvanaki2025sldb} & System-level SoC integration/configuration & No & No \\
ProtocolLLM~\cite{sheth2025protocolllm} & Protocol-specific (SPI/I\textsuperscript{2}C/UART/AXI) generation & No & No \\
\textbf{GateTruth (this work)} & Generation + agentic repair; audits external benchmarks & \textbf{Yes} & \textbf{Yes} \\
\bottomrule
\end{tabular}
\end{table}

NVIDIA's CVDP~\cite{pinckney2025cvdp} takes a substantially broader approach, spanning code generation, code comprehension, and both agentic and non-agentic evaluation across 783 problems in 13 task categories, with an Apache-2.0 harness and a public dataset whose non-code content is CC~BY~4.0 and whose code is Apache-2.0. CVDP is also the foundation of the Si2 (Silicon Integration Initiative) LLM Benchmarking Coalition, an industry effort co-chaired by NVIDIA's Nathaniel Pinckney and Synopsys's Ramesh Narayanaswamy, which extends CVDP with additional problem categories and operates a public results leaderboard (LBC-bench, launched 2026). As detailed in Section~\ref{sec:audit-method}, CVDP's public release deliberately withholds reference solutions to mitigate evaluation contamination; this is a sound design choice for CVDP's own scoring purpose but leaves the benchmark's own testbench rigor unauditable via mutation testing from the public release alone, a limitation we report rather than route around.

Research engaging with post-synthesis metrics has predominantly approached the problem from predictive modeling rather than generation scoring. MetRex~\cite{abdelatty2025metrex} compiled a comprehensive dataset of Verilog designs correlated with post-synthesis area, delay, and static power to evaluate whether models can infer these metrics from source code, while related frameworks like RocketPPA~\cite{abdollahi2025rocketppa} train models for code-level PPA prediction. Although these works advance metric estimation, they do not evaluate the physical viability of generated designs executed through an actual synthesis flow, nor do they address testbench rigor. GateTruth's own reference suite enforces a rigidly deterministic pipeline---a single content-addressed container image, uniform liberty and clock targets per task, a fixed Yosys-plus-OpenSTA path, and SHA-256 content hashes over canonical-JSON results---so that PPA comparisons produced by it are reproducible directly from the recorded digest.

In summary, GateTruth is, to our knowledge, the first tool to apply mutation-testing certification to existing public RTL-generation benchmarks and report the result, and the first to combine that audit methodology with an openly published, independently certified reference implementation used to validate it before external application. It uniquely combines: (a) a deterministic, seeded mutation-audit protocol applicable to any RTL benchmark with a golden reference and a runnable testbench; (b) an explicit, non-silent policy for benchmarks or designs that cannot be audited (\texttt{unsupported} status, or a documented structural finding, rather than a misleading score); (c) an openly published 68-task, dual-track reference suite in which the same protocol is certified against every one of the suite's 60 Track A testbenches before external use; and (d) a fully deterministic, reproducible RTL-to-PPA synthesis-and-timing flow underlying that reference suite's independent, secondary contribution as a model/agent benchmark.

\appendix

\section{Run-to-Run Variance}
\label{app:variance}

The v1.0 leaderboard reports a single run per model. Because model generation is not
deterministic (and the reasoning models do not accept a fixed temperature), we quantify
run-to-run spread by re-running the three Anthropic models over all 60 Track A tasks three
times each in a single fixed environment (pinned image, identical scoring). The v1.0 headline
numbers in Section~\ref{sec:results} are unchanged; Table~\ref{tab:variance} is an added
robustness measurement.

\begin{table}[h]
\centering
\caption{Track A run-to-run variance over three runs per model (GateTruth Score, 0--100). Standard deviation is the sample standard deviation (Bessel-corrected, $n-1$ denominator).}
\label{tab:variance}
\small
\begin{tabular}{lrrrr}
\toprule
Model & mean (n=3) & std & range & v1.0 (1 run) \\
\midrule
claude-opus-4-8   & 47.01 & 4.05 & 7.85 & 46.91 \\
claude-sonnet-4-6 & 47.89 & 3.90 & 7.78 & 45.03 \\
claude-haiku-4-5  & 35.05 & 0.98 & 1.89 & 33.58 \\
\bottomrule
\end{tabular}
\end{table}

Two observations follow. First, run-to-run standard deviation is substantial for the frontier
models---roughly 4 points, with a spread approaching 8 points across three runs---and much
smaller for Haiku (0.98). A single-run leaderboard is therefore a noisy estimator at the top of the
table; with only $n=3$ per model, these standard deviations are themselves imprecise, and we report
them as a directional signal rather than a tight interval. Second, and as a direct consequence,
Opus and Sonnet are statistically indistinguishable on Track A ($47.01 \pm 4.05$ vs $47.89 \pm
3.90$): their v1.0 single-run ordering lies well inside the noise, and across three runs Sonnet's
mean in fact edges ahead. The models' reliable separation appears in the agentic-repair track
(Section~\ref{sec:trackb}), not in static generation. We therefore recommend that future versions
report mean $\pm$ standard deviation over at least three runs as the default leaderboard cell, and
publish the individual per-run scores alongside the summary statistics so readers are not dependent
on our aggregation choices.

\bibliographystyle{plain}
\bibliography{references}

\end{document}
